% PCSF manuscript for Geophysics (SEG), Software and Algorithms section.
% Location: paper/manuscript of the pcsf repository. Class files: SEG
% template (S. Fomel, LPPL 1.3). Figures: python update_figures.py.
%
% Build:  pdflatex paper && bibtex paper && pdflatex paper && pdflatex paper
% Reproduce every number, table and figure (pcsf repository):
%   python paper/fetch_data.py && python paper/run_all.py
%
% \reviewtrue (default) builds the double-anonymous review version
% required by GEOPHYSICS: no acknowledgments, anonymised code link.
% \reviewfalse builds the accepted version.
%
% Red [PENDING: ...] markers: information only the author can supply.
% Bibliography entries marked VERIFY in pcsf.bib must be checked.
%
% Review mode: the default class option is "blind" (author block
% hidden). GEOPHYSICS uses double-anonymous review (instructions to
% authors, "Organizing your paper").
%
% The "manuscript" option follows SEG submission rules: double spacing,
% tables and figures collected at the end. For a reading copy with
% figures in the text, compile with \documentclass[paper]{geophysics}.
% When compiling elsewhere (e.g. Overleaf), upload the Fig/ and
% listings/ folders with paper.tex, pcsf.bib and the class files.

\documentclass[manuscript]{geophysics}

\usepackage{amsmath,amssymb}
\usepackage{xcolor}
\usepackage{booktabs}
\usepackage{array}
\usepackage{listings}
\usepackage{algorithm}
\usepackage{algpseudocode}
\usepackage{url}

\newif\ifreview
\reviewtrue
\DeclareRobustCommand{\pending}[1]{\textcolor{red}{\textbf{[PENDING:} #1\textbf{]}}}
\newcommand{\ohmm}{\ensuremath{\Omega\,\mathrm{m}}}
\newcommand{\code}[1]{\texttt{#1}}
\newcommand{\R}{\mathbb{R}}
% Geometry kinds (stored in a file as grid2d, grid3d,
% mesh_unstructured and multiline; Table tbl:geometry).
\newcommand{\Gtwo}{\mathsf{G}_{2}}
\newcommand{\Gthree}{\mathsf{G}_{3}}
\newcommand{\Mtri}{\mathsf{T}}
\newcommand{\Lset}{\mathsf{L}}
\newcolumntype{L}[1]{>{\raggedright\arraybackslash}p{#1}}
\algrenewcommand\algorithmicrequire{\textbf{Input:}}
\algrenewcommand\algorithmicensure{\textbf{Output:}}

\lstset{basicstyle=\ttfamily\scriptsize, breaklines=true,
        columns=fullflexible, frame=single}

% Figures live in Fig/ (the class default is the current folder).
% The class prefers Fig/<name>.pdf (vector) and falls back to .png.
\renewcommand{\figdir}{Fig}

\begin{document}

\title{A geometry-aware exchange format for electromagnetic
inversion models}

\renewcommand{\thefootnote}{\fnsymbol{footnote}}

\ms{GEO-XXXX}

\address{
\footnotemark[1]Laboratoire des Sciences G\'eographiques, du G\'enie
Civil et des G\'eosciences, \\
Institut National Polytechnique F\'elix Houphou\"et-Boigny, \\
BP 1093, Yamoussoukro, C\^ote d'Ivoire. \\
E-mail: etanoyau@gmail.com}
\author{Kouao Laurent Kouadio\footnotemark[1]}

\footer{PCSF inversion-model format}
\lefthead{Kouadio}
\righthead{Exchange format for EM inversions}

\maketitle

% =====================================================================
\begin{abstract}
Electromagnetic surveys have exchange standards for time series and
transfer functions but not for inversion models: each inversion code
writes its own format, and readers easily misread it. On five real
models, such misreadings caused errors of up to 4.5 decades for a
wrong logarithm base, 3.1 decades for an unstated depth datum, 2.7
decades for a mesh offset and 0.7 decades for a swapped axis order. We
present the Portable Common Subsurface Format (PCSF), which stores
isotropic resistivity on its native two-dimensional,
three-dimensional, triangular-mesh or multi-profile geometry. Each
file also keeps the source values and their encoding, the spatial
reference, and a record of how the model was obtained: data, errors,
solver, regularization, misfit, convergence and depth of
investigation. We derive and test the frame transformations, sampling
rules and uncertainty conversions needed to compare models of
different geometry. Three readers written without the reference
implementation read twelve files from Occam2D, ModEM, MARE2DEM and
synthetic predictors bit for bit. Three models of the EarthScope Earth
Model Collaboration passed through PCSF without loss; one archive
file, however, could not be reproduced from its solver model, and
none records its misfit or licence. Applied to 2D and 3D inversions of
the Baohuashan porphyry Cu--Mo survey, the comparison registered the
models to within 4.2~m. Within the depth of investigation, the models
agree near the surface but differ by 0.80 decades on average, mainly
where the 2D models contain resistive bodies that the 3D model does
not support.
\end{abstract}

% =====================================================================
\section{INTRODUCTION}

The resistivity model is the main product of a magnetotelluric (MT) or
controlled-source electromagnetic (EM) survey. It is the object that
geologists interpret, that later studies compare with seismic and
drilling results, and that new inversions use as a starting point or a
constraint. Outside the group that produced it, the model is useful
only if others can read it correctly; for an inversion result, this is
the practical meaning of the FAIR principles \citep{Wilkinson2016} and
of reproducible computational research \citep{Sandve2013}.

The data that enter an inversion have reached this point. SEG-Y serves
seismic data \citep{Hagelund2017}. For EM and MT, the SEG EDI standard
\citep{Wight1988} and its successor EMTF XML \citep{Kelbert2020}
describe transfer functions, and MTH5 \citep{Peacock2022MTH5} describes
time series. \citet{Kelbert2020} showed how much of the value of transfer
functions had been lost to ambiguous orientation, missing coordinates
and unstated sign conventions, and how a format with explicit metadata
recovered it. The resistivity model, the next link in the same chain
(Figure~\ref{fig:schematics}a), has had no comparable effort.

Inversion codes write models in formats designed around their solvers.
Occam2D \citep{deGrootHedlin1990} writes $\log_{10}$ resistivity on a
rectilinear mesh whose horizontal origin lies at the outer edge of the
lateral padding. ModEM \citep{Kelbert2014,Egbert2012} writes
natural-log resistivity on a 3D grid with the north axis first, places
the grid origin in an optional trailer, and fills the cells above the
topography with air. MARE2DEM \citep{Key2016} writes linear resistivity
per region and keeps the triangular mesh in separate files. Learned
predictors \citep{Rodriguez2023,Wang2023Prior,Corseri2024} follow no
convention at all. None of these files states which data were inverted,
how well they were fitted, or how deep they constrain the model. These
are the questions a reviewer asks first.

General-purpose containers solve only part of the problem. HDF5
\citep{Folk2011} and netCDF \citep{Rew1990} store typed arrays, and
VTK \citep{Schroeder2006} stores meshes for display, but none of them
states what the arrays of an EM model mean. Domain efforts cover parts
of that meaning. The IRIS Earth Model Collaboration distributes
gridded Earth models, MT models among them, as netCDF files
\citep{Trabant2012}. The European Plate Observing System includes a
service for MT data and conductivity models, and national programs in
the United States and Australia publish MT models from their own
archives \citep{Kelbert2020}. The UBC-GIF mesh and model files are a
common convention for tensor and octree meshes, and SimPEG
\citep{Cockett2015} and MTpy \citep{Krieger2014} read and write several
solver formats. Table~\ref{tbl:related} compares these options with
the requirements we set out below. Each covers some requirements, but
none was designed to carry a model from the common inversion codes
with the information needed to interpret it.

This paper describes the Portable Common Subsurface Format (PCSF) and
its supporting algorithms. We measure the ambiguities met when reading
the native files of three inversion codes, derive from them the
information a model file must contain, and define a data model with
binary and text encodings that holds it. We then give the frame
transformations, sampling rules and uncertainty conversions needed to
compare models of different geometry, test them for fidelity,
accuracy, independent readability and cost, and state which
conversions lose information. The comparison of 2D and 3D inversions
of one audio-magnetotelluric (AMT) survey, a standard question in MT
interpretation \citep{Ledo2005}, requires every one of these steps and
closes the paper.

\plot{schematics}{width=\textwidth}
{The format and its reference implementation. (a) The
electromagnetic data chain: time series and transfer functions have
community exchange standards (blue), whereas the resistivity model
produced by an inversion has none; PCSF fills this position (orange).
(b) Logical data model, with the symbols of
equation~\ref{eq:model}: a model has one geometry of a declared kind
($\Gtwo$, $\Gthree$, $\Mtri$ or $\Lset$), a canonical resistivity
field with its source values, their encoding and an optional
uncertainty, a station table, optional topography and an inversion
record; the same model is written losslessly to the HDF5 (PCSF) or
the text (PCSM) encoding. (c) Reference implementation: one adapter
per source converts solver output to the common model; the model is
written to either encoding and read by the sampling and comparison
functions, the conformance suite, viewers and independent readers.}

\tabl{related}{Existing representations compared with the requirements
R1--R6 defined in the text. $+$ supported; $\circ$ partly supported or
by convention only; $-$ not supported.
\pending{VERIFY each cell against the current documentation of EMC
netCDF, UBC-GIF, SimPEG/discretize and VTK.}}{
\centering\small
\begin{tabular}{@{}lcccccc@{}}
\toprule
Representation & R1 & R2 & R3 & R4 & R5 & R6 \\
\midrule
Native solver files (Occam2D, ModEM, MARE2DEM) & $\circ$ & $+$ & $\circ$ & $-$ & $-$ & $-$ \\
HDF5 or netCDF without conventions & $-$ & $\circ$ & $-$ & $-$ & $+$ & $-$ \\
netCDF with CF conventions \citep{Hassell2017} & $+$ & $\circ$ & $+$ & $\circ$ & $+$ & $+$ \\
IRIS EMC netCDF models & $+$ & $\circ$ & $+$ & $\circ$ & $+$ & $\circ$ \\
UBC-GIF mesh and model files & $\circ$ & $\circ$ & $\circ$ & $-$ & $\circ$ & $-$ \\
VTK \citep{Sullivan2019} & $-$ & $+$ & $-$ & $-$ & $\circ$ & $-$ \\
PCSF (this paper) & $+$ & $+$ & $+$ & $+$ & $+$ & $+$ \\
\bottomrule
\end{tabular}
}

% =====================================================================
\section{AMBIGUITIES IN INVERSION-MODEL FILES}

The models used in this paper come from three inversion codes and two
field surveys (Table~\ref{tbl:datasets}). Reading them correctly
required knowledge that none of the files states. We describe the
ambiguities in the order in which a reader meets them and measure, for
each, the error that a plausible wrong reading causes. All errors are
root-mean-square differences of $\log_{10}\rho$, in decades, between
the correct and the wrong reading of the same file
(Table~\ref{tbl:ambiguities}).

The most basic ambiguity is the scale of the stored numbers (A1 and
A2 in Table~\ref{tbl:ambiguities}). ModEM writes
the natural logarithm of resistivity and states it in a header keyword
that many readers ignore; Occam2D writes $\log_{10}$ and states it
nowhere. Reading the Baohuashan ModEM model as $\log_{10}$ values
displaces its 197\,737 earth cells by 4.4 decades on average, and
reading an Occam2D model as natural logarithms displaces it by 2.3
decades. These errors are large enough to be noticed, but only by a
reader who already knows what resistivities to expect.

The order of the horizontal axes is harder to detect (A3). ModEM
orders its grid north first, so that its $x$ axis points north and its $y$
axis east, the reverse of the usual map convention. A reader who
assumes east first places each station column in a different part of
the model. Below the 122 stations of the Baohuashan survey, from the
surface to 1~km depth, this error is 0.72 decades on average and
reaches 2.5 decades. Unlike an encoding error, it produces a plausible
model and is therefore unlikely to be noticed.

The vertical datum is stated even less explicitly (A4). In the
Baohuashan ModEM files,
station depths are measured from a model top at 224~m above sea level,
a value that appears only in a comment line of the data file. A reader
who starts each profile at the top of the model, instead of at the
station, reads the first tens of meters from the air and shifts the
rest of the profile; the error is 3.1 decades on average and reaches
11.8 decades. The air cells, which carry $10^{11}$ to
$10^{13}$~\ohmm\ in this model, pose a related problem (A5): a reader
who interpolates across the surface without removing them contaminates
the upper 100~m by 4.5 decades on average and up to 11.7 decades.

Occam2D raises the same question for the horizontal origin (A6):
its mesh coordinate starts at the outer edge of the lateral padding,
not at the first station. The offset is 0.7~km for the
Tongkeng mesh and 12.7~km for the Baohuashan L18 mesh, and reading the
mesh coordinate as station chainage displaces the two models by 0.8
and 2.7 decades. The Tongkeng run files also contain no coordinates
at all; the station positions and elevations had to be recovered from
a separate survey file.

A file may also hold another model than the one intended (A7).
MARE2DEM writes one resistivity file per iteration, and an example file
distributed with an earlier version of pyCSAMT contained the starting
model, written by a reader that has since been corrected. It passed
every self-consistency test because it was only compared with
itself. Its
6538 free regions differ from those of the final iteration by 0.77
decades on average and up to 2.4 decades.

These cases share three features. The information needed to read the
model correctly existed, but outside the model file or in a form that
software does not read. A wrong reading produced a model that looked
plausible in most cases. And the errors are of the same size as the
geological contrasts the models are meant to resolve. A format for
inversion models must therefore state this information explicitly, in
a form that software reads.

\tabl{ambiguities}{Ambiguities met in the native files and the error
caused by a wrong reading, as root-mean-square, mean absolute and
maximum absolute difference of $\log_{10}\rho$ (decades).}{
\centering\small
\begin{tabular}{@{}lL{5.4cm}L{3.0cm}ccc@{}}
\toprule
 & Wrong reading & Model & RMS & Mean & Max \\
\midrule
A1 & natural-log values read as $\log_{10}$ & ModEM, Baohuashan & 4.45 & 4.37 & 5.82 \\
A2 & $\log_{10}$ values read as natural log & Occam2D, L18 & 2.47 & 2.32 & 4.31 \\
A3 & north-first axes read as east-first & ModEM, Baohuashan & 0.72 & 0.58 & 2.54 \\
A4 & profiles started at the model top (datum) & ModEM, Baohuashan & 3.10 & 1.26 & 11.77 \\
A5 & air cells kept (0--100 m below surface) & ModEM, Baohuashan & 4.47 & 2.37 & 11.68 \\
A6 & mesh coordinate read as chainage (0.7 km) & Occam2D, Tongkeng & 0.79 & 0.62 & 3.79 \\
A6 & mesh coordinate read as chainage (12.7 km) & Occam2D, L18 & 2.67 & 2.28 & 5.54 \\
A7 & starting model read for the final one & MARE2DEM & 0.77 & 0.60 & 2.38 \\
\bottomrule
\end{tabular}
}

% =====================================================================
\section{CRITICAL INFORMATION FOR AN INVERSION MODEL}

The ambiguities above concern how to read the numbers; a second class
of information, just as often missing, concerns how to judge them.
\citet{Kelbert2020} defined the critical metadata of a transfer
function as the information without which it cannot be interpreted or
reused. For an inversion model, this information falls into seven
groups, which PCSF stores as an inversion record
(Table~\ref{tbl:record}).

The data group records what was inverted, from the components (for
example TE and TM apparent resistivity and phase, or the full
impedance) to the errors and the sign convention of the time
dependence; two models of the same survey can differ simply because
they fit different data. The forward and regularization groups record
the dimensionality and mesh of the forward problem and the type, final
trade-off parameter and starting model of the regularization, which
together select one of the many models that fit the data. The misfit
group records the final and target root-mean-square (RMS) misfit, the
number of iterations and how the run ended. We distinguish runs that
reached their target misfit, runs stopped by an iteration or time
limit, and runs that stalled, because a model from a run that stalled
at a misfit of 8 does not carry the weight of one that reached a
misfit of 1. The reliability group records the depth to which the data
constrain the model and, optionally, uncertainty or sensitivity arrays.
The solver and provenance (producing code, creation time and native
values, described in the next section) and the citation and licence
under which the model may be reused complete the record.

The inversion record has eight required and fifteen recommended
fields (Table~\ref{tbl:record}). The reference implementation checks
which are missing, so that a writer or an archive can decide whether a
file is complete enough to share. For the field models of this paper,
every field except the licence could be filled from the run files: the
Occam2D startup, data, log and configuration files, the ModEM data,
control and log files, and the MARE2DEM data and resistivity headers.
The check also flags real gaps. The Tongkeng run files do not document
the data errors, so one required field is missing, and the synthetic
test predictions, which are not inversions, lack seven of the eight
required fields (Table~\ref{tbl:completeness}).

\tabl{record}{The inversion record. R, required; r, recommended.}{
\centering\small
\begin{tabular}{@{}L{2.5cm}L{8.6cm}L{3.6cm}@{}}
\toprule
Group & Fields & Example (Baohuashan ModEM) \\
\midrule
Solver & \code{solver} (R), \code{solver\_version} (r) & ModEM 3D MT, NLCG \\
Data & \code{components} (R), \code{frequency\_range\_hz} (R), \code{n\_stations} (R),
  \code{errors} (R), \code{n\_frequencies}, \code{sign\_convention}, \code{source} (r) &
  $Z_{xy}$, $Z_{yx}$; 10.2--10\,400 Hz; 122; 7--31\% of $|Z|$; $e^{+i\omega t}$ \\
Forward & \code{dimensionality}, \code{mesh} (r) & 3D; $91 \times 54 \times 44$ \\
Regularization & \code{type} (R), \code{trade\_off}, \code{start\_model} (r) &
  model covariance; $1.2 \times 10^{-6}$; prior model \\
Misfit & \code{final\_rms} (R), \code{termination} (R), \code{target\_rms},
  \code{n\_iterations} (r) & 1.95; stalled; 1.20; 113 \\
Reliability & \code{method}, \code{depth\_of\_investigation\_m} (r) &
  Niblett--Bostick; per station \\
Citation & \code{authors}, \code{year}, \code{licence} (r); \code{title}, \code{doi} &
  \pending{licence} \\
\bottomrule
\end{tabular}
}

\tabl{completeness}{Completeness of the inversion records of the
twelve files (present / total).}{
\centering\small
\begin{tabular}{@{}lcccc@{}}
\toprule
File & Required & Recommended & Termination & Final RMS \\
\midrule
Tongkeng, Occam2D & 7 / 8 & 9 / 15 & target reached & 1.00 \\
Baohuashan L18--L34, Occam2D (5 files) & 8 / 8 & 14 / 15 & time limit (3), stalled (2) & 7.58--10.37 \\
Baohuashan multiline & 8 / 8 & 14 / 15 & per line & per line \\
Baohuashan, ModEM & 8 / 8 & 14 / 15 & stalled & 1.95 \\
Baohuashan sample run, ModEM & 8 / 8 & 10 / 15 & stalled & 3.06 \\
MARE2DEM demonstration & 8 / 8 & 7 / 15 & target reached & 1.00 \\
Synthetic predictions (2 files) & 1 / 8 & 1 / 15 & -- & -- \\
\bottomrule
\end{tabular}
}

% =====================================================================
\section{DATA MODEL AND ENCODINGS}

From the ambiguities and the critical information we derive six
requirements. The physical quantity, unit and scale of every value
must be stated in the file (R1). The model must be stored on the
discretization the solver used, without resampling (R2). Station
positions, topography, the coordinate reference system, the axis
convention and the depth datum must be stored with the model (R3). The
source values and their encoding, the producing code and the inversion
record must travel with it (R4). A reader must be able to interpret the
file with a general-purpose library and the specification alone, and a
text form must allow inspection and line-by-line comparison (R5).
Files must carry a version, so that readers detect files they cannot
interpret (R6). Figure~\ref{fig:schematics}b shows the logical data
model that meets these requirements.

\subsection{Formal definition}

A PCSF model is the tuple
\begin{equation}
M = \left(\kappa,\ \mathcal{G},\ \boldsymbol{\rho},\ \mathbf{r},\ e,\
\mathbf{U},\ \mathcal{T},\ \mathcal{H},\ \mathcal{I},\ v\right),
\qquad \kappa \in \mathcal{K},
\label{eq:model}
\end{equation}
where $\mathcal{K} = \{\Gtwo, \Gthree, \Mtri, \Lset\}$ is the set of
geometry kinds (a two-dimensional grid, a three-dimensional grid, a
triangular mesh and a set of profiles), $\mathcal{G}$ the geometry, $\boldsymbol{\rho}$ the canonical
field, $\mathbf{r}$ and $e$ the source values and their encoding,
$\mathbf{U}$ an optional uncertainty expressed in the encoding $e$,
$\mathcal{T}$ the station table, $\mathcal{H}$ the topography,
$\mathcal{I}$ the inversion record and $v$ the format version. The
geometry defines cells $\Omega_c$, $c = 0, \ldots, n_c - 1$, that
partition the model domain $\Omega = \bigcup_c \Omega_c$, and every
field is a vector indexed by the cells,
$\boldsymbol{\rho} \in \left((0, \infty) \cup
\{\mathrm{NaN}\}\right)^{n_c}$, where NaN marks cells without a value,
such as air or padding.

For $\kappa = \Gtwo$ and $\Gthree$, strictly monotonic
nodes $x_0, \ldots, x_{n_x}$, $y_0, \ldots, y_{n_y}$ and
$z_0, \ldots, z_{n_z}$ (depth positive downward) define the cells
\begin{equation}
\Omega_{ik} = [x_k, x_{k+1}] \times [z_i, z_{i+1}],
\qquad
\Omega_{ijk} = [x_k, x_{k+1}] \times [y_j, y_{j+1}] \times [z_i, z_{i+1}],
\label{eq:cells}
\end{equation}
with centers $\bar{x}_k = (x_k + x_{k+1})/2$ (and likewise
$\bar{y}_j$, $\bar{z}_i$), stored in row-major order with depth first:
$c = i\,n_x + k$ in 2D and $c = (i\,n_y + j)\,n_x + k$ in 3D. For
$\kappa = \Mtri$, nodes
$\mathbf{v}_0, \ldots, \mathbf{v}_{n_v - 1} \in \R^2$ and triangles
$(a_c, b_c, d_c)$ define
$\Omega_c = \operatorname{conv}\{\mathbf{v}_{a_c}, \mathbf{v}_{b_c},
\mathbf{v}_{d_c}\}$, and a region map
$\zeta: \{0, \ldots, n_c - 1\} \to \{0, \ldots, n_r - 1\}$ relates a
per-region table to the per-triangle field by
$\rho_c = \rho^{\mathrm{reg}}_{\zeta(c)}$. Region ids are therefore
0-based indices into the table; MARE2DEM numbers its regions from 1,
and its adapter converts the numbering once, on conversion, so that a
reader never has to guess the base of the index. For
$\kappa = \Lset$, the geometry is a set
$\{(\ell, \mathcal{G}_\ell, \boldsymbol{\rho}_\ell, y_\ell,
o_\ell)\}$ of $\Gtwo$ models with cross-line offsets $y_\ell$
and labels $o_\ell \in \{\text{measured}, \text{assigned}\}$, and the
model has no field of its own. A model is conformant when the
predicates of Table~\ref{tbl:rules} hold; the reference
implementation reports each of them.

\tabl{rules}{Conformance rules. A rule applies when the field it
names is present; NaN cells are exempt from C2 and C5.}{
\centering\small
\begin{tabular}{@{}lL{4.4cm}L{8.6cm}@{}}
\toprule
Rule & Name & Predicate \\
\midrule
C1 & geometry kind & $\kappa \in \mathcal{K}$ \\
C2 & positivity & $\rho_c > 0$ for every cell with a value \\
C3 & shape & $\boldsymbol{\rho}$ has one value per cell of $\mathcal{G}$
  (for meshes, per triangle or per region) \\
C4 & encoding & $e \in \mathcal{E}$ when $\mathbf{r}$ is stored \\
C5 & source values & $\max_c \lvert C_e(r_c) - \rho_c \rvert / \rho_c
  \le 10^{-12}$ \\
C6 & multiline root & a $\Lset$ model has no field of its own \\
C7 & derived volume & a volume interpolated between lines is labeled
  synthesized \\
C8 & geographic coordinates & station longitude and latitude are both
  present or both absent \\
C9 & uncertainty & $\mathbf{U}$ and the sensitivity have the shape of
  $\boldsymbol{\rho}$ \\
C10 & region ids & $0 \le \zeta(c) < n_r$ for every triangle of a
  $\Mtri$ model \\
\bottomrule
\end{tabular}
}

\subsection{Encodings and the canonical field}

The canonical field is resistivity $\rho$ in linear ohm-meters (R1).
Adapters convert the source values $r$, with declared encoding
$e \in \{\log_{10}, \ln, \mathrm{lin}, \log_{10}\sigma, \sigma\}$, by
\begin{equation}
\rho = C_e(r) =
\begin{cases}
10^{\,r}, & e = \log_{10} \quad \text{(Occam2D)},\\
\exp(r),  & e = \ln       \quad \text{(ModEM)},\\
r,        & e = \mathrm{lin} \quad \text{(MARE2DEM)},\\
10^{-r},  & e = \log_{10}\sigma \quad \text{(EMC archive)},\\
1/r,      & e = \sigma,
\end{cases}
\label{eq:canonical}
\end{equation}
where the last two encodings store conductivity $\sigma$ in S/m, as
community archives do. Adapters store $r$ and $e$ unchanged next to
$\rho$ (R4), so that a reader can use either the common scale or the
exact values the solver wrote. Comparisons in this paper
use the log-resistivity $m = \log_{10}\rho$.

Each $C_e$ is a bijection from its domain $D_e$ ($\R$ for the
logarithmic encodings, $(0, \infty)$ otherwise) onto $(0, \infty)$,
with inverse $C_e^{-1}$ and derivative
\begin{equation}
\frac{\mathrm{d}\rho}{\mathrm{d}r} =
\rho\ln 10,\quad \rho,\quad 1,\quad -\rho\ln 10,\quad -\rho^{2}
\qquad \text{for } e = \log_{10},\ \ln,\ \mathrm{lin},\
\log_{10}\sigma,\ \sigma .
\label{eq:jacobian}
\end{equation}
A change of encoding $e \to e'$, $r' = C_{e'}^{-1}\!\left(C_e(r)\right)$,
therefore has the derivative
$\mathrm{d}r'/\mathrm{d}r =
(\mathrm{d}\rho/\mathrm{d}r)_{e} / (\mathrm{d}\rho/\mathrm{d}r')_{e'}$,
a constant between logarithmic encodings. Rule C5
(Table~\ref{tbl:rules}) checks that the stored pair $(r, e)$
reproduces $\rho$.

A reader must read the declared kind $\kappa$ before it interprets
any array, because array shape alone does not identify the geometry
(Table~\ref{tbl:geometry}). The kinds also fix conventions that the
native files leave open: a $\Gtwo$ profile uses the station chainage
as its horizontal coordinate, with the Occam2D padding offset removed
(ambiguity A6); a $\Gthree$ volume keeps the ModEM axis order and
stores its origin, rotation and number of air layers; and a volume
interpolated between the profiles of a $\Lset$ model may be stored but
is always labeled as synthesized (rule C7).

\tabl{geometry}{The four geometry kinds, with the identifier that a
file stores for each.}{
\centering\small
\begin{tabular}{@{}L{2.6cm}L{4.0cm}L{3.0cm}L{5.0cm}@{}}
\toprule
Kind $\kappa$ (stored as) & Typical source & Resistivity shape & Spatial reference \\
\midrule
$\Gtwo$, 2D grid (\texttt{grid2d}) & Occam2D; 2D predictors & $(n_z, n_x)$ &
  station chainage; optional origin and azimuth \\
$\Gthree$, 3D grid (\texttt{grid3d}) & ModEM; 3D predictors & $(n_z, n_y, n_x)$ &
  grid origin, rotation, air layers \\
$\Mtri$, triangular mesh (\texttt{mesh\_unstructured}) & MARE2DEM; graph predictors &
  per triangle (per region, per node optional) &
  node coordinates; mesh plane \\
$\Lset$, set of profiles (\texttt{multiline}) & set of 2D inversions & one $(n_z, n_x)$ per line &
  cross-line offset per line \\
\bottomrule
\end{tabular}
}

The station table links the model frame to the Earth (R3). It stores,
for every station, its name, its coordinates in the model frame
(including its depth below the model datum), its elevation, an
optional line identifier, and its longitude and latitude in WGS84. A
root attribute stores the coordinate reference system as a string
accepted by the PROJ library, and topography is stored per station or
as a gridded surface. With this table, the frame transformations of
the next section can be estimated from the file itself, which resolves
ambiguities A3 and A4 without outside knowledge.

Every file records the producing code, the writing software, the
creation time and the resistivity unit. An optional history group
stores per-iteration quantities such as the misfit and the trade-off
parameter, and the inversion record is stored as JSON text in the
metadata. For a model produced directly by a learned predictor, an
additional record stores the architecture, framework, checkpoint
identifier and its SHA-256 hash, training-data reference,
hyperparameters and random seed. The hash confirms that a checkpoint
file is the one recorded; it does not show that the checkpoint
produced the model, which requires running the predictor again.

The binary encoding (PCSF) is a single HDF5 file
(Appendix~\ref{app:layout}). The text encoding (PCSM) writes the same
model as keyword blocks, defined by a grammar in extended
Backus--Naur form in the specification of the format, with
floating-point values written as the shortest decimal strings that
convert back to the same IEEE 754 numbers \citep{Steele1990}, so the
two encodings convert into each other without loss. PCSM relates to
PCSF as CDL text relates to netCDF. Each file carries a semantic
version (R6): a reader rejects an unknown major version and warns
about a newer minor version. The version described here is 0.1.0.

% =====================================================================
\section{FRAME TRANSFORMATIONS, SAMPLING AND UNCERTAINTY}

Comparing two models requires placing them in one frame, evaluating
them at the same points, carrying their uncertainties across
encodings, and restricting the comparison to where the data constrain
them. This section derives each step. Table~\ref{tbl:notation}
summarizes the notation.

\tabl{notation}{Notation.}{
\centering\small
\begin{tabular}{@{}L{2.6cm}L{11.8cm}@{}}
\toprule
Symbol & Meaning \\
\midrule
$M,\ \mathfrak{M}$ & a model (equation~\ref{eq:model}) and the set of conformant models \\
$\kappa,\ \mathcal{K}$ & geometry kind and the set of kinds \\
$\Gtwo,\ \Gthree,\ \Mtri,\ \Lset$ & kinds: 2D grid, 3D grid, triangular mesh, set of profiles \\
$\mathcal{G},\ \Omega_c,\ n_c$ & geometry, its cells and their number \\
$\mathcal{E}$ & set of encodings \\
$r,\ e$ & source values and their encoding \\
$\rho,\ m$ & canonical resistivity ($\Omega$\,m) and $m = \log_{10}\rho$ \\
$C_e$ & conversion from source values to $\rho$ (equation~\ref{eq:canonical}) \\
$\mathbf{J},\ \mathbf{C}_r,\ \mathbf{C}_m$ & Jacobian of a change of encoding; covariance of $r$ and of $m$ \\
$\mathbf{u},\ \mathbf{g}$ & projected (easting, northing) and model-frame horizontal coordinates \\
$\mathbf{P},\ \mathbf{R}(\theta),\ \mathbf{S}$ & axis permutation, rotation and scale of a model frame \\
$z_{\mathrm{top}},\ z_s$ & elevation of the model datum and depth of station $s$ below it \\
$\Pi_a,\ \Pi_g$ & arithmetic and geometric node-to-cell projection \\
$\mathcal{S}_M(\mathbf{p})$ & sampled log-resistivity of model $M$ at point $\mathbf{p}$ \\
$w_c(\mathbf{p}),\ \tau$ & interpolation weight of cell $c$ and axis fraction (equation~\ref{eq:weights}) \\
$\mathcal{P},\ D$ & target points of a comparison and interpolated depth of investigation \\
$W,\ R$ & writer and reader of an encoding ($h$, HDF5; $t$, text) \\
$\sigma_m,\ \sigma_\rho$ & standard deviation of $m$ and of $\rho$ \\
$h_{\mathrm{NB}},\ D_s$ & Niblett--Bostick depth and depth of investigation of station $s$ \\
$\Delta,\ \varepsilon,\ \beta,\ R$ & difference, root-mean-square misfit, bias and correlation \\
$h,\ q$ & cell size and convergence order, $\varepsilon \propto h^{q}$ \\
\bottomrule
\end{tabular}
}

\subsection{Frame transformations}

Let $\mathbf{u} = (E, N)^{T}$ be the projected easting and northing of
a point, for example in UTM coordinates, and $\mathbf{g}$ its
horizontal coordinates in a model frame. For a grid whose origin is
$\mathbf{u}_0$, whose horizontal axes are rotated by $\theta$
clockwise from grid north and scaled by $\mathbf{S}$, and whose axes
are stored in the order given by a permutation $\mathbf{P}$, the two
are related by
\begin{equation}
\mathbf{g} = \mathbf{A}\,\mathbf{u} + \mathbf{b},
\qquad
\mathbf{A} = \mathbf{R}(\theta)\,\mathbf{S}\,\mathbf{P},
\qquad
\mathbf{b} = -\mathbf{A}\,\mathbf{u}_0 ,
\label{eq:frame}
\end{equation}
with $\mathbf{P} = \mathbf{I}$ for an east-first grid and
$\mathbf{P} = \bigl(\begin{smallmatrix}0 & 1\\ 1 & 0\end{smallmatrix}\bigr)$
for a north-first grid such as ModEM's. The determinant of $\mathbf{A}$
is negative for a north-first grid, so a reader detects ambiguity A3
from the numbers alone. Vertically, a point at depth $d$
below the surface at station $s$ lies at model depth
\begin{equation}
z = z_s + d,
\qquad
z_s = z_{\mathrm{top}} - H_s ,
\label{eq:datum}
\end{equation}
where $H_s$ is the station elevation and $z_{\mathrm{top}}$ the
elevation of the model datum. Storing $z_s$ in the station table makes
$z_{\mathrm{top}}$ unnecessary for comparisons (ambiguity A4). For an
Occam2D profile, the horizontal coordinate is $x = x_{\mathrm{mesh}} -
W_{\mathrm{pad}}$, where $W_{\mathrm{pad}}$ is the width of the left
padding (ambiguity A6); the adapter applies this shift once, on
conversion.

When the station table stores both $\mathbf{u}_i$ and $\mathbf{g}_i$
for $n_s$ stations, $\mathbf{A}$ and $\mathbf{b}$ are estimated by least
squares,
\begin{equation}
(\hat{\mathbf{A}}, \hat{\mathbf{b}}) = \arg\min_{\mathbf{A}, \mathbf{b}}
\sum_{i=1}^{n_s} \left\lVert \mathbf{A}\mathbf{u}_i + \mathbf{b} - \mathbf{g}_i \right\rVert^2 ,
\label{eq:register}
\end{equation}
and the singular-value decomposition of $\hat{\mathbf{A}}\mathbf{P}$
returns $\theta$ and $\mathbf{S}$. The residual measures how well the
frame is described by equation~\ref{eq:frame}.

\subsection{Node-to-cell projection}

Graph networks predict one value per mesh node, whereas the canonical
mesh field is defined per triangle. For a triangle $t$ with vertex set
$V(t)$, the arithmetic and geometric projections are
\begin{equation}
\Pi_a:\ \rho(t) = \frac{1}{|V(t)|}\sum_{v \in V(t)} \rho_{\mathrm{node}}(v),
\qquad
\Pi_g:\ m(t) = \frac{1}{|V(t)|}\sum_{v \in V(t)} m_{\mathrm{node}}(v).
\label{eq:node}
\end{equation}
$\Pi_g$ is the geometric mean of the node resistivities. Because
resistivity varies over orders of magnitude, $\Pi_a$ is dominated by
the most resistive vertex. Both projections are available, and the
choice is recorded in the file.

\subsection{Sampling}

The sampling operator evaluates the log-resistivity of a model $M$ at
a point $\mathbf{p}$ of its frame,
\begin{equation}
\mathcal{S}_M(\mathbf{p}) =
\begin{cases}
\displaystyle\sum_{c} w_c(\mathbf{p})\, m_c, & \kappa \in \{\Gtwo, \Gthree\},\\[6pt]
m_{t(\mathbf{p})}, & \kappa = \Mtri,
\end{cases}
\label{eq:sample}
\end{equation}
and is undefined outside the model. On grids, $m_c$ are cell-center
values and $w_c$ the bi- or trilinear weights of the surrounding
cells; between the outermost cell center and the outermost cell edge,
the edge value is used. On meshes, $t(\mathbf{p})$ is the triangle
that contains $\mathbf{p}$, found by point location, so that contrasts
across region boundaries are kept. Interpolation is applied to $m$,
not to $\rho$, so that relative changes count equally in resistive
and conductive parts of a model. In a model with topography, the air cells above
each column take the value of the first earth cell below them before
sampling, so that interpolation near the surface involves earth values
only (ambiguity A5).

On one axis with cell centers $\xi_0 < \cdots < \xi_{n-1}$, a
coordinate $q$ between the outer cell edges is first clipped to
$\tilde{q} = \min(\max(q, \xi_0), \xi_{n-1})$ and then receives the
weights
\begin{equation}
w_i = 1 - \tau,\quad w_{i+1} = \tau,
\qquad
\tau = \frac{\tilde{q} - \xi_i}{\xi_{i+1} - \xi_i},
\qquad
i = \max\{\, j \le n-2 : \xi_j \le \tilde{q} \,\},
\label{eq:weights}
\end{equation}
and the weight of a grid cell is the product of its axis weights,
$w_c(\mathbf{p}) = \prod_a w^{(a)}_{c_a}(p_a)$, with four cells in 2D
and eight in 3D. The weights are non-negative and sum to one, so
$\mathcal{S}_M$ is a convex combination of cell values: it never
leaves the range of the values it interpolates, and it reproduces any
field that is affine within the bracketing cells. For a field $m$
with continuous second derivatives, sampled at the cell centers, the
error of tensor-product linear interpolation is bounded by
\begin{equation}
\left\lvert \mathcal{S}_M(\mathbf{p}) - m(\mathbf{p}) \right\rvert
\le \frac{1}{8} \sum_{a} h_a^{2}
\max \left\lvert \frac{\partial^{2} m}{\partial a^{2}} \right\rvert ,
\label{eq:interp-error}
\end{equation}
where $h_a$ is the spacing of the cell centers along axis $a$ and the
maxima are taken over the bracketing cells. Cell averages differ from
center values by $O(h^2)$, so the order is unchanged: sampling is of
second order in $h$ where $m$ is smooth and of first order across
discontinuities, such as the interfaces of the model of Test~2.

\subsection{Uncertainty under a change of encoding}

A stored uncertainty is meaningful only in the encoding it refers
to. Between logarithmic encodings the conversion is exact,
\begin{equation}
\sigma_{\ln\rho} = \ln(10)\,\sigma_m ,
\label{eq:sigma-log}
\end{equation}
because the two logarithms differ by a constant factor. Toward linear
resistivity, the first-order (delta-method) approximation is
\begin{equation}
\sigma_\rho \approx \rho\,\ln(10)\,\sigma_m ,
\qquad
\frac{\sigma_\rho}{\rho_{\mathrm{med}}} =
\sqrt{\left(e^{s^2} - 1\right) e^{s^2}},\quad s = \ln(10)\,\sigma_m ,
\label{eq:sigma-lin}
\end{equation}
where the second expression is the exact standard deviation of a
log-normal resistivity with median $\rho_{\mathrm{med}}$. The
approximation underestimates the exact value by less than 1\% for
$\sigma_m \le 0.05$ decades, by 3.9\% at 0.1, by 15\% at 0.2 and by 64\%
at 0.5 decades (Table~\ref{tbl:delta}). PCSF therefore stores
uncertainty in the encoding of the source values, and the reference
implementation converts it with equations~\ref{eq:sigma-log} and
\ref{eq:sigma-lin}; for $\sigma_m$ above about 0.1 decades, the
interval $\rho \cdot 10^{\pm\sigma_m}$ should be reported instead of a
symmetric linear standard deviation.

The same rules apply to a covariance. For source values with
covariance $\mathbf{C}_r$, the covariance in another encoding is, to
first order,
\begin{equation}
\mathbf{C}_{r'} \approx \mathbf{J}\,\mathbf{C}_r\,\mathbf{J}^{T},
\qquad
\mathbf{J} = \operatorname{diag}\!\left(
\frac{\mathrm{d}r'_c}{\mathrm{d}r_c}\right),
\label{eq:cov}
\end{equation}
with the derivatives of equation~\ref{eq:jacobian}; the relation is
exact between logarithmic encodings, where $\mathbf{J}$ is constant.
Because the sampling operator is linear in $m$, the variance of a
sampled value follows exactly from the weights of
equation~\ref{eq:weights},
\begin{equation}
\sigma_{\mathcal{S}}^{2}(\mathbf{p}) =
\mathbf{w}(\mathbf{p})^{T}\,\mathbf{C}_m\,\mathbf{w}(\mathbf{p}),
\qquad
\sum_c w_c^{2}\,\sigma_c^{2}
\ \le\ \sigma_{\mathcal{S}}^{2}(\mathbf{p})
\ \le\ \Big(\sum_c w_c\,\sigma_c\Big)^{2},
\label{eq:interp-var}
\end{equation}
where $\sigma_c^2$ are the diagonal elements of $\mathbf{C}_m$. The
lower bound is reached for independent cells and requires
non-negative correlations between them; the upper bound, reached for
fully correlated cells, holds for any covariance because the weights
are non-negative and
$\lvert (\mathbf{C}_m)_{cd} \rvert \le \sigma_c\,\sigma_d$.

\tabl{delta}{Relative standard deviation of linear resistivity from a
log-uncertainty $\sigma_m$: first-order approximation
(equation~\ref{eq:sigma-lin}, left) against the exact log-normal
value.}{
\centering\small
\begin{tabular}{@{}lcccccc@{}}
\toprule
$\sigma_m$ (decades) & 0.01 & 0.05 & 0.1 & 0.2 & 0.3 & 0.5 \\
\midrule
First order & 0.023 & 0.115 & 0.230 & 0.461 & 0.691 & 1.151 \\
Exact & 0.023 & 0.116 & 0.240 & 0.540 & 0.993 & 3.225 \\
Error (\%) & $-0.04$ & $-1.0$ & $-3.9$ & $-14.8$ & $-30.4$ & $-64.3$ \\
\bottomrule
\end{tabular}
}

\subsection{Depth of investigation}

A difference between two models is meaningful only where both are
constrained by data. We use the Niblett--Bostick depth
\citep{Jones1983,Spies1989}, computed for each station from the
rotation-invariant (determinant) impedance
$Z_{\det} = (Z_{xx}Z_{yy} - Z_{xy}Z_{yx})^{1/2}$:
\begin{equation}
\rho_a(T) = 0.2\,T\,\lvert Z_{\det}(T)\rvert^{2},
\qquad
h_{\mathrm{NB}}(T) = \sqrt{\frac{\rho_a(T)\,T}{2\pi\mu_0}} \approx 356\sqrt{\rho_a(T)\,T}\ \mathrm{m},
\qquad
D_s = \max_T h_{\mathrm{NB}}(T),
\label{eq:doi}
\end{equation}
with $Z$ in mV/km/nT, $T$ the period in seconds and $\mu_0$ the
magnetic permeability of free space. $D_s$ is a simple, data-based
estimate; sensitivity-based estimates \citep{Oldenburg1999} can be
stored in the same field when the inversion provides them. The values
are stored per station in the reliability group of the inversion
record and interpolated along a profile when models are compared.

\subsection{Comparison}

Two models $A$ and $B$ are compared at a common set of $N$ target
points where both are defined and that lie above the depth of
investigation (Algorithm~\ref{alg:compare}),
\begin{equation}
\mathcal{P} = \left\{ \mathbf{p}_k :\ \mathbf{p}_k \in \Omega_A \cap
\Omega_B,\ \ d_k \le D(\mathbf{p}_k) \right\}_{k=1}^{N},
\label{eq:targets}
\end{equation}
where $d_k$ is the depth of $\mathbf{p}_k$ below the surface and $D$
the depth of investigation interpolated between the stations. The
difference and its summary statistics are
\begin{equation}
\Delta_k = \mathcal{S}_A(\mathbf{p}_k) - \mathcal{S}_B(\mathbf{p}_k),
\qquad
\varepsilon = \Big[\frac{1}{N}\sum_{k}\Delta_k^{2}\Big]^{1/2},
\qquad
\beta = \frac{1}{N}\sum_{k}\Delta_k,
\qquad
R = \operatorname{corr}\big(\mathcal{S}_A, \mathcal{S}_B\big),
\label{eq:stats}
\end{equation}
where $\varepsilon$ is in decades, $\beta$ shows whether one model is
systematically more resistive, and $R$, the Pearson correlation of the
two sampled fields, shows whether they place structures in the same
positions regardless of their absolute level. The same statistics,
with $B$ replaced by an exact field, measure the accuracy of sampling.

\begin{algorithm}[htbp]
\caption{Registration and comparison of two PCSF models.}
\label{alg:compare}
\begin{algorithmic}[1]
\Require models $A$, $B$; projection to a common coordinate system;
  target points $\{(\mathbf{u}_k, d_k)\}_{k=1}^{N}$ (projected position,
  depth below surface); air threshold $\rho_{\mathrm{air}}$
\Ensure $\Delta$, $\varepsilon$, $\beta$, $R$
\For{$X \in \{A, B\}$}
  \State estimate $(\mathbf{A}_X, \mathbf{b}_X)$ from the station table of $X$
    \Comment{equation~\ref{eq:register}}
  \State replace each cell with $\rho_X > \rho_{\mathrm{air}}$ by the first earth cell below it
  \For{$k = 1, \ldots, N$}
    \State $\mathbf{g} \gets \mathbf{A}_X \mathbf{u}_k + \mathbf{b}_X$;\quad
      $z \gets z_s(\mathbf{g}) + d_k$ \Comment{equations~\ref{eq:frame}--\ref{eq:datum}}
    \State $s_{X,k} \gets \mathcal{S}_X(\mathbf{g}, z)$ \Comment{equation~\ref{eq:sample}}
  \EndFor
\EndFor
\State $\mathcal{K} \gets \{k : s_{A,k}, s_{B,k}$ defined and $d_k \le D(\mathbf{u}_k)\}$
  \Comment{equation~\ref{eq:doi}}
\State compute $\Delta$, $\varepsilon$, $\beta$, $R$ over $\mathcal{K}$
  \Comment{equation~\ref{eq:stats}}
\end{algorithmic}
\end{algorithm}

% =====================================================================
\section{IMPLEMENTATION AND EXTENSIBILITY}

The reference implementation is pcsf, an open-source Python package
published separately from any inversion or processing library
(Figure~\ref{fig:schematics}c). Its core contains the data classes,
the HDF5 and text readers and writers, the conformance rules, the
inversion record with its completeness check, and the functions
\code{sample\_log10}, \code{compare\_log10} and
\code{log\_uncertainty}, which implement
equations~\ref{eq:sample}--\ref{eq:stats}; it needs only Python 3.9 or
newer, NumPy and h5py 3.8 or newer \citep{Collette2013}. One adapter
per source (\code{occam2d\_to\_pcsf}, \code{modem\_to\_pcsf},
\code{mare2dem\_to\_pcsf} and \code{emc\_to\_pcsf}) reads the native
files with parsers that use only NumPy, so that each solver format is
parsed in one place and the format depends on no inversion package.
Every adapter follows the same procedure
(Algorithm~\ref{alg:conversion}): it converts the source values,
projects node values when needed, attaches the geometry, the station
table and the inversion record, and checks shapes and positivity. The
package also provides JSON Schemas \citep{Pezoa2016} of the inversion
record and of the provenance of learned models, a conformance suite of
reference files with their expected checksums, a command-line tool and
a viewer. pyCSAMT \citep{Kouadio2022} uses pcsf as its format core and
adds the assembly of profiles into $\Lset$ models.

Adding a solver requires one adapter that fills the same fields; no
change to the format or to the readers is needed, because readers
dispatch on the declared geometry kind, not on the producing code.
Adapters, encodings and companion formats of other packages register
through Python entry points, without changes to pcsf.
Adding a new geometry kind or a new physical quantity, such as
anisotropic resistivity or chargeability, is a minor version change
under the versioning rules: older readers ignore what they do not
know and still read the rest of the file.

\begin{algorithm}[htbp]
\caption{Conversion of solver output to a PCSF model.}
\label{alg:conversion}
\begin{algorithmic}[1]
\Require source values $r$ with encoding $e$; geometry kind $\kappa$ and
  geometry $G$; projection $\Pi \in \{\Pi_a, \Pi_g\}$; stations $S$;
  inversion record $I$; optional topography $T$, uncertainty $U$, history $H$
\Ensure validated model $M$
\State $\rho \gets C_e(r)$ \Comment{equation~\ref{eq:canonical}}
\If{$\kappa = \Mtri$ and $r$ is given per node}
  \State $\rho_{\mathrm{node}} \gets \rho$;\quad $\rho \gets \Pi(\rho_{\mathrm{node}})$
    \Comment{equation~\ref{eq:node}; record $\Pi$}
\EndIf
\State $M \gets (\kappa, G, \rho, r, e, S, T, U, H, I)$
\State check positivity of $\rho$, its shape against $G$, and the completeness of $I$
\State \Return $M$
\end{algorithmic}
\end{algorithm}

% =====================================================================
\section{NUMERICAL TESTS}

Table~\ref{tbl:datasets} lists the models. The tests check the format
and the algorithms; they do not check whether a solver result is
geologically correct. The scripts that produce every number, table and
figure are archived with the code.

\tabl{datasets}{Models used in the tests and examples.}{
\centering\small
\begin{tabular}{@{}L{4.3cm}L{1.9cm}L{2.9cm}L{6.0cm}@{}}
\toprule
Model & Code & Kind & Size and run \\
\midrule
Tongkeng CSAMT profile \citep{Kouadio2020} & Occam2D & $\Gtwo$ &
  47 stations; TM mode, 10--10\,000 Hz; $31 \times 576$ cells; RMS 1.00 \\
Baohuashan lines L18, L22, L26, L30, L34 \citep{Kouabena2025} & Occam2D &
  $\Gtwo$, $\Lset$ & 25 stations per line; TE and TM, 42--10\,400 Hz;
  RMS 7.6--10.4 \\
Baohuashan 3D model & ModEM & $\Gthree$ & 122 stations; $Z_{xy}$, $Z_{yx}$,
  10--10\,400 Hz; $44 \times 54 \times 91$ cells; RMS 1.95 \\
Baohuashan 3D sample run & ModEM & $\Gthree$ & 125 stations; full $Z$;
  590\,400 cells; RMS 3.06 \\
MARE2DEM demonstration & MARE2DEM & $\Mtri$ &
  synthetic MT data; 6774 triangles, 6540 regions; RMS 1.00 \\
Synthetic 2D prediction & generic & $\Gtwo$ & $14 \times 24$ cells \\
Synthetic graph prediction & generic & $\Mtri$ &
  63 nodes, 96 triangles \\
Synthetic accuracy model & generic & $\Gtwo$, $\Mtri$ & exact field; $h$ = 25--400~m \\
EMC MCR-MT-2016, THO-MT-2021, WUS-MT-2021 \citep{Bedrosian2016,Bedrosian2021,Murphy2021} &
  ModEM; EMC netCDF & $\Gthree$ & 1.2, 4.7 and 4.0 million archive values;
  solver files $150 \times 150 \times 50$ to $220 \times 210 \times 85$ \\
\bottomrule
\end{tabular}
}

\subsection{Test 1: Fidelity to the solver output}

For the MARE2DEM model, the per-region table in the PCSF file was
compared with the resistivity column of the final iteration file,
parsed by the test script itself: all 6540 values are equal and in region order,
and the mesh triangulated from the model description with Triangle
\citep{Shewchuk1996} contains all 6540 region identifiers. For both
ModEM models, the source values in the PCSF file are the same 590\,400
and 216\,216 numbers as the model block of the solver file, parsed
independently. For all ten files that store source values, the
canonical array equals $C_e(r)$ exactly, and for the Tongkeng model
the source array equals the $31 \times 576$ model returned by the
Occam2D result reader. It was this comparison with the solver files,
not the round trips of the format, that exposed the starting model in
the earlier MARE2DEM example (ambiguity A7).

\subsection{Test 2: Accuracy of sampling and projection}

The sampling and projection algorithms change the support of the
values, so their accuracy was measured against an exact field
(Figure~\ref{fig:test2_sampling}a). The model $m^\star(x, z)$ on a
$4 \times 2$~km section has a 100~\ohmm\ layer over a 1000~\ohmm\
half-space at 400~m depth, a conductor of $-1.5$ decades dipping
$35^\circ$ and a resistive block of $+0.8$ decades. It was discretized
on quality triangular meshes and on regular grids with cell size $h$
from 400 to 25~m, and evaluated on a 5~m grid of 321\,201 points
(Table~\ref{tbl:test2}).

Linear interpolation of $m$ on grids was the most accurate method at
every cell size below 400~m and had the highest convergence order
($q = 1.56$); at 25~m its error was 0.0028 decades, against 0.0043
for interpolation of $\rho$ and 0.022 for the nearest cell. On meshes,
point-location sampling converged at first order ($q = 0.94$), as
expected for a piecewise-constant field. The geometric projection
$\Pi_g$ was more accurate than the arithmetic projection $\Pi_a$ at
every cell size, by 12\% at 400~m and 14\% at 200~m, and both approached
the discretization floor of the mesh at small $h$; $\Pi_g$ should
therefore be the default for graph predictors.

\tabl{test2}{RMS error $\varepsilon$ (decades of $\Omega$\,m) against
the exact model, and fitted convergence order $q$.}{
\centering\small
\begin{tabular}{@{}lcccc@{}}
\toprule
Method & $h$ = 400 m & $h$ = 100 m & $h$ = 25 m & $q$ \\
\midrule
Mesh, exact cell value (point location) & 0.178 & 0.054 & 0.0132 & 0.94 \\
Mesh, $\Pi_a$ (arithmetic node mean) & 0.252 & 0.064 & 0.0135 & 1.06 \\
Mesh, $\Pi_g$ (geometric node mean) & 0.223 & 0.057 & 0.0133 & 1.02 \\
Grid, linear in $m = \log_{10}\rho$ & 0.213 & 0.032 & 0.0028 & 1.56 \\
Grid, linear in $\rho$ & 0.266 & 0.052 & 0.0043 & 1.49 \\
Grid, nearest cell & 0.252 & 0.084 & 0.0223 & 0.88 \\
\bottomrule
\end{tabular}
}

\plot{test2_sampling}{width=\textwidth}
{Accuracy test. (a) Exact log-resistivity model. (b) Triangular mesh
with $h = 200$~m and cell values from the geometric node projection
$\Pi_g$. (c) RMS error $\varepsilon$ against cell size for mesh
sampling and node projections (solid) and grid resampling (dashed);
fitted orders in parentheses.}

\subsection{Test 3: Reading without the reference implementation}

To test whether the specification alone is sufficient (R5), each of the
twelve files was read by three readers that do not import the
reference implementation: a
Python reader that uses only h5py and the specification; the HDF5
command-line tool \code{h5dump}, which exports datasets as raw binary
through the HDF5 C library; and a parser written from the PCSM
grammar of the specification, applied to the text file. All 851\,390
resistivity values were identical, bit for bit, to those of the
reference reader for all three readers, and all 75 conformance checks
of the h5py reader passed (Table~\ref{tbl:readers}).

\tabl{readers}{Independent reading of the twelve files. ``Exact''
means bit-identical IEEE 754 values.}{
\centering\small
\begin{tabular}{@{}lrcccc@{}}
\toprule
Kind (files) & Values & h5py & h5dump & PCSM parser & Conformance \\
\midrule
$\Gtwo$ (7) & 28\,048 & exact & exact & exact & 42 / 42 \\
$\Gthree$ (2) & 806\,616 & exact & exact & exact & 12 / 12 \\
$\Mtri$ (2) & 6\,870 & exact & exact & exact & 11 / 11 \\
$\Lset$ (1) & 9\,856 & exact & exact & exact & 10 / 10 \\
\midrule
Total (12) & 851\,390 & exact & exact & exact & 75 / 75 \\
\bottomrule
\end{tabular}
}

\subsection{Test 4: Storage and input/output}

Each of four models was written and read five times in each
representation, and the median time is reported
(Figure~\ref{fig:test4_storage}). The netCDF-4 and VTK files hold only
the arrays a viewer needs and serve as lower bounds; they carry
neither the source values nor the station table nor the inversion
record.
PCSF was smaller than the native model files for three of the four
models: 0.10 against 0.16~MB for Tongkeng, 0.20 against 1.17~MB for
MARE2DEM and 4.1 against 7.7~MB for the ModEM sample; for the five
Occam2D lines it was 0.16 against 0.15~MB. For the ModEM volume, PCSF
was about twice the size of the minimal netCDF file (2.0~MB) because it
stores both $\rho$ and $r$. Absolute times depend on the computer and
its load, but their ratios were stable between repeated runs: reading
the 590\,400-cell volume took about 12 times longer from PCSM than
from PCSF (0.59 against 0.05~s in the run reported here), and
compressed PCSM, the smallest text form (3.9~MB), took about 25 times
longer to write than PCSF. PCSM is therefore suited to profiles,
inspection and version control, and PCSF to volumes.

\plot{test4_storage}{width=\textwidth}
{Storage test for the Tongkeng ($\Gtwo$), Baohuashan
($\Lset$), MARE2DEM ($\Mtri$) and ModEM sample ($\Gthree$)
models. (a) File size. (b) Write time. (c) Read time; median of five
repetitions. netCDF-4 and VTK files contain only coordinates and
resistivity.}

% =====================================================================
\section{ROUND TRIP THROUGH A COMMUNITY ARCHIVE}

A format for models must also work with the archives that already
distribute them. The IRIS (now EarthScope) Earth Model Collaboration
(EMC) holds 3D MT models as netCDF files and, for some of them, the
original solver files. We downloaded three such models: MCR-MT-2016
\citep{Bedrosian2016,EMC_MCR}, THO-MT-2021
\citep{Bedrosian2021,EMC_THO} and WUS-MT-2021
\citep{Murphy2021,EMC_WUS}. Each was computed with ModEM and is
distributed as a netCDF-4 file of $\log_{10}$ conductivity on a regular
0.1$^\circ$ grid of latitude and longitude, accompanied by the ModEM
model file and a short text that states the transverse Mercator
projection of the solver grid. The archive files thus use a sixth
convention, $\log_{10}\sigma$, in addition to the three of the solvers,
which PCSF now accepts as a declared encoding
(equation~\ref{eq:canonical}). The files themselves are not uniform:
the value array is named \code{log10sigma} in one and
\code{log\_10sigma} in the other two, and only the latter declare a
fill value.

We first tested whether PCSF carries an archive file without loss. Each
netCDF file was converted to PCSF, with its values as source values in
the $\log_{10}\sigma$ encoding, its latitude and longitude as model
coordinates in the WGS84 reference system, and every global and
variable attribute stored with its type in the metadata. The PCSF file
was written to PCSM, read back, written to PCSF again and converted
back to netCDF. For all three models, the 9\,783\,850 values were
bit-identical to the original, including their single-precision type,
and so were the coordinates, the file format and all 159 attributes
(Table~\ref{tbl:emc}). The PCSF files were 3.3 times larger than the
netCDF files, because they store both the canonical and the source
values in double precision, and the PCSM files 17 to 21 times larger,
which confirms that the text encoding is not meant for volumes of
millions of cells.

We then compared each archive file with the solver file it was made
from. The ModEM model was converted to PCSF and sampled at the archive
grid points with equation~\ref{eq:sample}, after projecting latitude
and longitude with the stated projection. The text gives two readings
of the projection that differ by 1.0 to 1.8~km, the false northing and
the latitude of the grid center, and we tried both, with linear
interpolation and with nearest-cell lookup. For WUS-MT-2021, nearest-cell
lookup with the grid-center reading reproduced 94.5\% of the archive
values exactly ($\varepsilon = 0.063$ decades), and the remaining
differences lie on cell boundaries and along the central meridian,
where nearest-cell lookup is sensitive to sub-cell differences in the
projection (Figure~\ref{fig:emc_vs_solver}a--c). For THO-MT-2021, the
false-northing reading did better (77\% of values reproduced,
$\varepsilon = 0.086$) than the grid-center reading (68\%). The correct
reading of the projection therefore differs between two models from the
same author and archive. For MCR-MT-2016, no reading and no sampling rule
reproduced the archive: under all four combinations, $\varepsilon$ was
0.245 decades, fewer than 0.1\% of the values agreed, and 0.4\% differed
by more than one decade (Figure~\ref{fig:emc_vs_solver}d--f). The
archive file shows the same structures as the solver file, with
sharper contrasts, and was evidently made from a different version of
the model or by a different procedure, which neither file documents.

The archive versions are also smaller than the solver models. They omit
10 to 18\% of the horizontal cells of the solver grids, mainly the
padding, and for THO-MT-2021 and WUS-MT-2021 the six deepest of 85
layers. Their metadata are rich in citation information, with authors,
year, reference and repository identifiers, but the reference
identifier of THO-MT-2021 in the file differs from the one on the
product page, that of WUS-MT-2021 is empty, and none of the files has a
licence. Mapped onto the inversion record, each archive file fills one
of the eight required fields, the solver; the data, errors,
regularization, misfit and convergence are absent
(Table~\ref{tbl:emc}).

The archive format thus makes models discoverable and citable, and
PCSF carries it without loss, but the archive version is a derived
product whose derivation is not recorded; it cannot replace the solver
model when a result has to be checked.

\tabl{emc}{Round trip of three EMC archive models through PCSF and
comparison of each archive file with its ModEM solver file. Equal
cells: fraction of archive values reproduced to $10^{-4}$ decades by
the best combination of projection reading and sampling rule.}{
\centering\small
\begin{tabular}{@{}lccc@{}}
\toprule
 & MCR-MT-2016 & THO-MT-2021 & WUS-MT-2021 \\
\midrule
Archive grid (depth $\times$ lat $\times$ lon) & $50 \times 129 \times 180$ & $79 \times 169 \times 350$ & $79 \times 200 \times 250$ \\
Solver grid (N $\times$ E $\times$ depth) & $150 \times 150 \times 50$ & $190 \times 220 \times 85$ & $220 \times 210 \times 85$ \\
Round trip: values, coordinates & bit-equal & bit-equal & bit-equal \\
Round trip: attributes & 57 / 57 & 51 / 51 & 51 / 51 \\
Size: netCDF / PCSF / PCSM (MB) & 3.1 / 10.3 / 53.7 & 10.5 / 34.7 / 220.8 & 10.0 / 33.2 / 186.7 \\
Best projection reading & none reproduces & false northing & grid center \\
Best sampling rule & none reproduces & nearest cell & nearest cell \\
Equal cells & $<$0.1\% & 77.0\% & 94.5\% \\
$\varepsilon$, best (decades) & 0.245 & 0.086 & 0.063 \\
$\varepsilon$, other reading, nearest cell (decades) & 0.245 & 0.101 & 0.097 \\
Solver horizontal cells inside archive & 81.7\% & 89.6\% & 85.4\% \\
Inversion record, required fields & 1 / 8 & 1 / 8 & 1 / 8 \\
\bottomrule
\end{tabular}
}

\plot{emc_vs_solver}{width=\textwidth}
{EMC archive files compared with their ModEM solver files. (a--c)
WUS-MT-2021 at 15.6~km depth: solver model resampled by nearest-cell
lookup with the grid-center projection, archive values, and their
difference; 94\% of the cells are equal. (d--f) MCR-MT-2016 at
14.8~km depth, solver model resampled by linear interpolation: the
archive file cannot be reproduced from the solver file under any
reading of the stated projection.}

% =====================================================================
\section{WHAT A CONVERSION KEEPS}

\citet{Kelbert2020} showed that the choice of a format decides which
scientific operations remain possible, and distinguished lossless from
lossy conversions between transfer-function formats. The same
distinction applies to inversion models, and our tests quantify it
(Table~\ref{tbl:lossy}). Let $\mathfrak{M}$ be the set of conformant
models and $W_h$, $W_t$ the writers and $R_h$, $R_t$ the readers of
the HDF5 and text encodings. The encodings are lossless in the sense
that
\begin{equation}
R_h(W_h(M)) = M,\qquad R_t(W_t(M)) = M,\qquad
R_t\big(W_t(R_h(W_h(M)))\big) = M
\qquad \text{for all } M \in \mathfrak{M},
\label{eq:lossless}
\end{equation}
with every stored array equal in IEEE 754 binary64 (Test~3). A
conversion $F$ is lossy when no map $G$ satisfies $G \circ F =
\mathrm{id}$ on $\mathfrak{M}$: regridding a mesh assigns one value to
every grid point inside a triangle, so the positions of region
boundaries within a grid step cannot be recovered, and the archive
resampling of the EMC replaces the native cells by a coarser support.

Conversion from a native solver file to PCSF keeps the source values
unchanged (Test~1) and the native geometry, and it is the step at which
the adapter adds what the native file lacks: the padding offset, the
axis convention, the station depths and the inversion record. The
other conversions of Table~\ref{tbl:lossy} keep the values but differ
in what they drop. A netCDF or VTK file of coordinates and resistivity
can be displayed, but without the source encoding, the station table
and the record it can no longer be checked against the data or
compared with another model. A native solver file written back from
PCSF loses the record, because the native formats have no place for
it. Regridding a mesh or interpolating a 3D volume between profiles
changes the support of the values, with the errors measured in
Test~2, and PCSF labels such a result as synthesized so that it is
never mistaken for an inversion result.

\tabl{lossy}{What each conversion keeps. L, lossless; lossy conversions
name what is lost.}{
\centering\small
\begin{tabular}{@{}L{3.6cm}L{2.2cm}L{8.6cm}@{}}
\toprule
From $\rightarrow$ to & Values & What is lost \\
\midrule
Native $\rightarrow$ PCSF & L & nothing; offset, axis order, datum and record are made explicit \\
PCSF $\leftrightarrow$ PCSM & L & nothing (bit-exact, Test~3) \\
PCSF $\rightarrow$ netCDF (arrays only) & L & source encoding, stations, inversion record \\
PCSF $\rightarrow$ VTK & L & source encoding, stations, record, coordinate reference \\
Mesh $\rightarrow$ grid; lines $\rightarrow$ 3D volume & lossy & support of the values; error per Test~2 \\
PCSF $\rightarrow$ native solver file & L & inversion record, station positions (Occam2D) \\
EMC netCDF $\leftrightarrow$ PCSF & L & nothing (bit-exact values and attributes) \\
Solver file $\rightarrow$ EMC netCDF & lossy & native grid, padding, deepest layers, source
  encoding, derivation; resampled to 0.1$^\circ$ \\
\bottomrule
\end{tabular}
}

% =====================================================================
\section{FIELD EXAMPLES}

\subsection{Tongkeng CSAMT profile}

Figure~\ref{fig:field_tongkeng} shows the Tongkeng Occam2D model
\citep{Kouadio2020} drawn from its PCSF file alone. The run files
contain no coordinates; the station elevations and positions were
added to the file from the survey station list by interpolation along
the surveyed path, and the plotting code uses only the geometry kind
and the station table.

\plot{field_tongkeng}{width=\textwidth}
{Occam2D model of the Tongkeng CSAMT profile (47 stations) drawn from
its PCSF file below the station topography, without vertical
exaggeration.}

\subsection{MARE2DEM demonstration model}

Figure~\ref{fig:field_mare2dem} shows the MARE2DEM model on its native
mesh and the same window sampled on a 20~m grid with
equation~\ref{eq:sample}. The sampled image keeps the triangle
boundaries of the native model because each point takes the value of
the triangle that contains it.

\plot{field_mare2dem}{width=\textwidth}
{MARE2DEM demonstration model, a marine survey. (a) Native
triangular mesh colored by resistivity below the seafloor (solid
line); the fixed seawater regions (0.3~\ohmm) are shown in blue and the
air is blank. Triangles mark the 13 seafloor receivers. (b) Detail of
the dashed window with triangle edges. (c) The same window sampled on
a 20~m grid by point location. Vertical exaggeration 0.8 in (a) and
1.3 in (b) and (c).}

\subsection{2D and 3D inversions of the Baohuashan survey}

The Baohuashan AMT survey consists of five parallel lines over a
porphyry Cu--Mo system \citep{Kouabena2025}. Each line was inverted
with Occam2D from the processed TE and TM apparent resistivity and
phase (27 frequencies, 42~Hz to 10.4~kHz, error floors of 10\% and
$2.86^\circ$), with a limit of 900~s per run. Three runs stopped at the
time limit and two stalled before the target misfit, with final RMS
misfits of 7.6 to 10.4. An independent ModEM 3D inversion of the
off-diagonal impedances of 122 stations of the same survey (19
frequencies, 10~Hz to 10.4~kHz) stalled in a local minimum at an RMS
misfit of 1.95 after 113 iterations.
\pending{Name the colleague who ran the ModEM inversion
(acknowledgments).}
These facts come from the inversion records of the files
(Table~\ref{tbl:completeness}). The five Occam2D models were written as
$\Gtwo$ files with the station longitude, latitude and elevation
from the EDI headers, and combined into a $\Lset$ file whose
offsets, computed from the station coordinates, are 0, 201, 401, 599
and 799~m (Figure~\ref{fig:field_baohuashan}).

\plot{field_baohuashan}{width=\textwidth}
{Baohuashan $\Lset$ model. (a) Station lines in UTM zone 50N.
(b) The five Occam2D sections at their measured cross-line offsets,
drawn in elevation below the station topography (solid lines), with
the stations (triangles); air is blank.}

The registration of the ModEM grid to UTM coordinates, estimated from
its 122 stations with equation~\ref{eq:register}, recovered the
north-first permutation (determinant $-1.001$), unit scales (1.0014 and
0.9996) and a rotation of $-0.03^\circ$, with an RMS residual of 4.2~m
and a maximum of 6.9~m, less than a tenth of the station spacing. The
18\,479 ModEM air cells were replaced as described above. The depth of
investigation from equation~\ref{eq:doi}, computed from the 27
processed frequencies and stored per station in the inversion records,
ranges from 340~m to 4.5~km over the stations, with a median of
1.9~km along the lines. Target points were
placed every 20~m along each line and every 10~m to 1.5~km below the
surface; 82\,601 of them (91\%) lie above the depth of investigation
and are used in the statistics.

Figure~\ref{fig:field_2d3d_sections} shows the two models along each
line and their difference, Figure~\ref{fig:field_2d3d_slices} shows
them at three depths, and Table~\ref{tbl:field} lists the statistics.
Within the depth of investigation, $\varepsilon = 1.29$ decades,
$\beta = -0.80$ and $R = 0.22$: the 3D model is less resistive and the
two models are weakly correlated. Without the depth limit the values
are almost the same ($\varepsilon = 1.35$, $\beta = -0.84$), so the
disagreement is not an artifact of comparing unconstrained depths. It
depends strongly on depth. In the upper 250~m, $\varepsilon = 0.81$
and $R = 0.59$, and the two models place the same near-surface
conductors. Between 500 and 1000~m, $\varepsilon = 1.59$ and
$\beta = -1.22$.

The disagreement follows the resistive bodies of the 2D models. Within
the depth of investigation, the fraction of 2D values above
$10^{5}$~\ohmm\ is zero on L26, 3\% on L22, 7\% on L18 and 32--33\% on
L30 and L34, and the bias grows from $0.04$ on L26 to $-1.26$ on L30.
On L18, part of the resistive block lies below the depth of
investigation, and excluding it reduces the bias from $-0.93$ to
$-0.65$. Values of $10^{5}$ to $10^{6}$~\ohmm\ are at the upper end of
the range of crystalline rocks and unusual for the altered rocks of a
porphyry system, and the 3D model, whose misfit is four to five times
lower, does not contain them. The comparison therefore identifies
these bodies as the parts of the 2D interpretation that need the most
caution: they may result from 3D effects that a 2D inversion cannot
represent \citep{Ledo2005} or from runs that did not reach their
target misfit, which the inversion records show to be the case for
all five lines. On L26, the line without such bodies, the two models
agree in level ($\beta = 0.04$, $\varepsilon = 0.51$).
\pending{Relate the shallow conductors common to both models and the
deep 3D conductor to the geology and drilling of
\citet{Kouabena2025}.}

\tabl{field}{Comparison of the Occam2D and ModEM models of the
Baohuashan survey within the depth of investigation
($\Delta = m_{3D} - m_{2D}$; 0--1.5~km).}{
\centering\small
\begin{tabular}{@{}lccccccc@{}}
\toprule
Line or depth & Occam2D RMS & DOI median (m) & 2D $> 10^5$ \ohmm & $\varepsilon$ & $\beta$ & $R$ & $N$ \\
\midrule
L18 & 7.80 & 1606 & 7\% & 1.33 & $-0.65$ & $-0.12$ & 14\,928 \\
L22 & 9.07 & 1734 & 3\% & 1.24 & $-0.98$ & 0.31 & 15\,788 \\
L26 & 7.58 & 1795 & 0\% & 0.51 & 0.04 & 0.35 & 17\,223 \\
L30 & 10.37 & 2532 & 32\% & 1.61 & $-1.26$ & 0.02 & 17\,369 \\
L34 & 7.68 & 2104 & 33\% & 1.47 & $-1.16$ & 0.34 & 17\,293 \\
All lines & & 1916 & 15\% & 1.29 & $-0.80$ & 0.22 & 82\,601 \\
\midrule
0--250 m & & & 4\% & 0.81 & $-0.28$ & 0.59 & 14\,950 \\
250--500 m & & & 18\% & 1.38 & $-0.80$ & 0.09 & 14\,900 \\
500--1000 m & & & 23\% & 1.59 & $-1.22$ & 0.07 & 27\,588 \\
1000--1500 m & & & 13\% & 1.08 & $-0.67$ & $-0.07$ & 25\,163 \\
\bottomrule
\end{tabular}
}

\plot{field_2d3d_sections}{width=\textwidth}
{Baohuashan lines L18 to L34, drawn in elevation below the station
topography (solid line); air is blank. (a) Occam2D models. (b) ModEM
model sampled along each line with Algorithm~\ref{alg:compare}, at the
same depths below the ground. (c) Difference $\Delta$; positive where
the 3D model is more resistive. The dashed line is the depth of
investigation, and the veiled area below it is excluded from the
statistics, which are given within each panel. Triangles mark
stations. Vertical exaggeration 0.8.}

\plot{field_2d3d_slices}{width=\textwidth}
{Baohuashan models at 100, 400 and 1000~m below the surface. In each
pair, the left panel shows the ModEM model (background) and the
Occam2D models along the lines (filled circles, every 100~m) on the
same color scale; the right panel shows the difference $\Delta$. Open
circles mark points below the depth of investigation.}

% =====================================================================
\section{DISCUSSION}

The ambiguities of the native files are not rare accidents. In five
models from three codes we met seven, and each caused errors of the
same size as the geological contrasts the models are meant to show.
Most of them produce plausible models, so they are unlikely to be
caught by inspection. PCSF removes them by construction: the encoding,
axis convention, datum, offset and iteration are either applied once
by the adapter or stated in the file, and the frame transformation
estimated from the station table turns a stated convention into a
checked one.

The inversion record addresses a different weakness. A model without
its data, misfit and convergence state cannot be weighed, and in our
field example these facts change the interpretation: all five 2D runs
stopped before their target misfit, the first explanation to
consider for the resistive bodies that the 3D model does not contain.
The record makes this visible in the file itself, and its completeness
check shows what a file still lacks, as for the Tongkeng errors and the
licences here. Because the record also carries the depth of
investigation, the comparison could show that the disagreement lies
within the depth the data resolve.

The algorithms carry measurable errors, and the tests set their
defaults: interpolation of log-resistivity on grids, point location on
meshes because it keeps the native boundaries, the geometric node
projection, and the first-order conversion of uncertainty only below
about 0.1 decades. For learned predictors, PCSF
stores the predicted model in a form that other groups can read and
compare with classical inversions; reproducing a prediction still
requires the checkpoint, the input data and the inference code, and
our tests of this path use synthetic arrays.

The present version has limitations. Only isotropic resistivity is
supported. Meshes are limited to triangles, so tetrahedral and octree
models are not yet represented. The frame transformation assumes an
affine relation between a model frame and projected coordinates, which
holds for the rectilinear and rotated grids of current codes. The
Niblett--Bostick depth is an approximate, data-based limit; when an
inversion provides sensitivities, a sensitivity-based depth of
investigation should be stored instead. The format has one reference
implementation; its metadata are described by JSON Schemas and its
rules by a conformance suite, but a second, independent
implementation of the writers is still needed before version 1.0.

A format is useful only if others adopt it. The specification and
the reference implementation are published separately from pyCSAMT,
with their own versions. Readers and
writers can be offered to MTpy and SimPEG, the inversion record can be
aligned with the metadata vocabulary of MTH5 \citep{Peacock2022MTH5}
and EMTF XML \citep{Kelbert2020}, and archives that already hold MT
models could accept PCSF files with their records. Because PCSF
already carries the archive files of the Earth Model Collaboration
without loss, storing each solver model with its record next to the
archive file would make the archive version traceable to the model it
was derived from. The SEG Technical Standards Committee, which
maintains EDI, would be a natural place to discuss a community
version.

% =====================================================================
\section{CONCLUSIONS}

Inversion models are exchanged in formats that leave their scale, axis
order, datum, origin and iteration to the reader's knowledge, and in
five real models from three codes the corresponding misreadings caused
errors of 0.7 to 4.5 decades. PCSF stores an EM inversion model on its
native geometry with these conventions resolved and keeps a
record of the data, errors, solver, regularization, misfit,
convergence and depth of investigation, in a binary and a text
encoding that convert into each other without loss and that three
independent readers read bit for bit. It also carries the netCDF files
of a community archive without loss, and comparing three of them with
their solver models showed that archive versions are derived products
whose derivation, misfit and licence are not recorded. With the frame
transformations,
sampling rules and uncertainty conversions given here, a 2D and a 3D
inversion of the Baohuashan survey were registered to within 4.2~m and
compared within their depth of investigation, which showed agreement
near the surface and located the deep resistive bodies of
non-converged 2D runs as the main difference.

% =====================================================================
% GEOPHYSICS: no acknowledgments section before acceptance
% (double-anonymous review).
\ifreview\else
\section{ACKNOWLEDGMENTS}

The author thanks Rujun Chen, leader of GSAI Geophysics Ltd, for
permission to use the Baohuashan data to test the format and for the
advancement of science, and the developers of h5py, HDF5 and the
Triangle mesh generator.
\pending{Name and thank the colleague who provided the ModEM 3D
inversion; funding; data providers for the Tongkeng survey.}
\pending{AI disclosure as required by SEG, for example: ``The author
used [tool, version] to draft and edit parts of the manuscript and the
analysis scripts. The author reviewed and revised all text and code
and is responsible for the content.'' VERIFY the required wording and
placement in the Ethical Guidelines for SEG Publications.}
\fi

\section{DATA AND MATERIALS AVAILABILITY}

The reference implementation, pcsf, is available at
\ifreview
\pending{anonymised review link (Zenodo restricted access or
anonymous.4open.science)}
\else
\url{https://github.com/earthai-tech/pcsf}
\fi
under the BSD 3-Clause licence; the repository also holds the
specification of the format, its JSON Schemas and its conformance
suite, and its source code accompanies the submission. The folder
\code{paper} of that repository reproduces every number, table and
figure of this paper: \code{run\_all.py} runs all analyses with pcsf
alone from the archived models and native solver files;
\code{run\_all.py --rebuild} first rebuilds the models from the solver
runs and EDI files with pyCSAMT \citep{Kouadio2022,Kouadio2026Zenodo};
and \code{check\_results.py} compares the output with the values
reported here.

The data archive \citep{PCSFdata} holds the models in PCSF and PCSM,
the native Occam2D, ModEM and MARE2DEM runs, the processed EDI files of
Baohuashan lines L18 and L22, and the three EMC models. The measured
data of the other Baohuashan lines are not public: the archive holds
the inversion models of all lines, but not the EDI files and Occam2D
data files of lines L26, L30 and L34 or the data files of the two
ModEM runs, which are available from the corresponding author on
reasonable request. Every result of the paper can be reproduced
without them; only the rebuild of those models needs them. The EMC
models are openly available from the EarthScope Earth Model
Collaboration \citep{EMC_MCR,EMC_THO,EMC_WUS} and are included in the
archive unchanged.
\pending{DOI of the data record \citep{PCSFdata}; agreement of the
author of the Baohuashan ModEM model to its release.}

% =====================================================================
\append[app:layout]{HDF5 layout}

Table~\ref{tbl:layout} lists the groups, datasets and attributes of a
PCSF file. Names are exact.

\tabl{layout}{Layout of a PCSF (HDF5) file. Optional entries are marked
with an asterisk.}{
\centering\footnotesize
\begin{tabular}{@{}L{4.4cm}L{2.4cm}L{8.2cm}@{}}
\toprule
Path & Type & Content \\
\midrule
\code{/} attributes & strings & \code{pcsf\_version}, \code{source\_backend},
  \code{created\_by}, \code{created\_at}, \code{resistivity\_unit} (always
  \code{ohm.m}), \code{crs}*, \code{description}* \\
\code{geometry/} & group & attribute \code{kind}; datasets of the kind
  (Table~\ref{tbl:geometry}) \\
\code{model/resistivity} & float array & canonical $\rho$ (absent for
  $\Lset$) \\
\code{model/resistivity\_native}* & float array & source values $r$;
  attribute \code{encoding} \\
\code{model/resistivity\_by\_region}* & float array & per-region table (meshes) \\
\code{model/resistivity\_by\_node}* & float array & per-node values (meshes) \\
\code{model/uncertainty}*, \code{sensitivity}* & float arrays & shape of
  \code{resistivity} \\
\code{stations/}* & group & \code{name}, \code{x}, \code{y}, \code{z},
  \code{line\_id}*, \code{lon}*, \code{lat}* (WGS84) \\
\code{topography/}* & group & attribute \code{kind}: \code{per\_station}
  or \code{raster} \\
\code{history/}* & group & per-iteration series, for example RMS misfit \\
\code{survey\_json}*, \code{metadata\_json}* & string & JSON text, including
  the inversion record (\code{inversion}), the node projection and the
  predictor record \\
\code{geometry/lines/<id>/} & group & $\Lset$ only: nested
  $\Gtwo$ geometry, \code{resistivity}, offset and its kind \\
\bottomrule
\end{tabular}
}

\append[app:record]{Inversion record of the Baohuashan ModEM model}

The inversion record is stored as JSON in \code{metadata\_json}. The
following listing is the record of the Baohuashan ModEM file, with the
per-station depths of investigation shortened for printing.

\lstinputlisting{listings/baohuashan_modem3d_record.json}

\newpage
\bibliographystyle{seg}
\bibliography{pcsf}

\end{document}
